\documentclass[12pt]{article}

\usepackage{bbm}
\usepackage{fancyhdr}
\usepackage{titling}
\usepackage{enumitem}
\usepackage{amsmath}
\usepackage{dsfont}
\usepackage{amsfonts}
\usepackage{amssymb}
\usepackage{amsthm}
\usepackage{geometry}
\usepackage{tikz}
\usepackage{algorithm}
\usepackage{algpseudocode}
\geometry{verbose,tmargin=1in,bmargin=1in,lmargin=1in,rmargin=1in}
\renewcommand{\vec}[1]{\mathbf{#1}}
\newcommand{\dec}{\text{dec}}
\newcommand{\divi}{\mathbin{\textnormal{div}}}
\title{Homework 1}
\author{Satvik Prasad}
\date{}
\newtheorem{lemma}{Lemma}[section]
\fancypagestyle{plain}{%
  \fancyhf{}
  \fancyhead[L]{\theauthor}
  \fancyhead[R]{\thedate}
}
\pagestyle{plain}
\makeatletter
\def\@maketitle{%
  \newpage
  \null
  \vskip 1em%
  \begin{center}%
    \let \footnote \thanks
    {\LARGE \@title \par}%
    \vskip 1em%
  \end{center}%
  \par
\vskip 1em}
\makeatother
\usepackage{cmbright}

\newcommand{\mat}[1]{\mathbf{#1}}

% ---------- Probability theory ----------
% Core operators (optional argument-free; use as \E{X}, \Var{X}, \Cov{X}{Y})
\DeclareMathOperator{\EOp}{\mathbb{E}}
\DeclareMathOperator{\VarOp}{Var}
\DeclareMathOperator{\CovOp}{Cov}
\DeclareMathOperator{\CorrOp}{Corr}

\renewcommand{\P}[1]{\mathbb{P}\left(#1\right)}        % \P{A}
\newcommand{\E}[1]{\EOp\left[#1\right]}                % \E{X}
\newcommand{\Var}[1]{\VarOp\left(#1\right)}            % \Var{X}
\newcommand{\Cov}[2]{\CovOp\left(#1,\,#2\right)}       % \Cov{X}{Y}
\newcommand{\Corr}[2]{\CorrOp\left(#1,\,#2\right)}     % \Corr{X}{Y}
\newcommand{\SD}[1]{\operatorname{SD}\left(#1\right)}  % \SD{X}

% Subscripted versions, e.g. \Esub{\theta}{X}, \Psub{\theta}{A}
\newcommand{\Esub}[2]{\EOp_{#1}\left[#2\right]}
\newcommand{\Psub}[2]{\mathbb{P}_{#1}\left(#2\right)}

% Conditioning, indicators, independence
\newcommand{\given}{\,\middle|\,}                      % \E{X \given Y}
\newcommand{\ind}[1]{\mathbbm{1}\left\{#1\right\}}     % \ind{X > 0}
\newcommand{\indep}{\perp\!\!\!\perp}                  % X \indep Y
\newcommand{\iid}{\overset{\text{iid}}{\sim}}          % X_i \iid \Normal(0,1)

% Convergence
\newcommand{\convp}{\xrightarrow{\;p\;}}               % in probability
\newcommand{\convd}{\xrightarrow{\;d\;}}               % in distribution
\newcommand{\convas}{\xrightarrow{\text{a.s.}}}        % almost surely
\newcommand{\convLp}[1]{\xrightarrow{L^{#1}}}          % in L^p

% Probability spaces
\newcommand{\F}{\mathcal{F}}                           % sigma-algebra
\newcommand{\B}{\mathcal{B}}                           % Borel sigma-algebra
\newcommand{\pspace}{(\Omega, \F, \mathbb{P})}

% Common distributions
\newcommand{\Normal}{\mathcal{N}}
\newcommand{\Bern}{\operatorname{Bernoulli}}
\newcommand{\Bin}{\operatorname{Binomial}}
\newcommand{\Pois}{\operatorname{Poisson}}
\newcommand{\Unif}{\operatorname{Uniform}}
\newcommand{\Exp}{\operatorname{Exp}}
\newcommand{\Geom}{\operatorname{Geometric}}
\newcommand{\Gam}{\operatorname{Gamma}}
\newcommand{\Beta}{\operatorname{Beta}}

% Number sets
\newcommand{\R}{\mathbb{R}}
\newcommand{\N}{\mathbb{N}}
\newcommand{\Z}{\mathbb{Z}}


\begin{document}
\subsection{Writing Proofs}
\begin{enumerate}
  \item Assume we have some surjective linear map $T : V \to W,$
    where $V$ is finite dimensional. Then, we know we can construct a
    basis for $V$, $\beta$, such that
    $$
    \beta = \{v_1, \ldots, v_n, u_1, \ldots, u_m\},
    $$
    and $v_1, \ldots, v_n$ is a basis for $\ker T,$ because $\ker T$
    is a subspace of $V.$ Then, let $U = \text{span} \{ u_1, \ldots, u_m\},$ and
    define $T|_U : U \to W$ such that
    $$
    T|_U(u) = Tu.
    $$

    Because $T$ is surjective, for all $w \in W,$ there exists some
    $u \in V$ such that $Tv = w.$ We can write $v = \lambda_1 v_1 +
    \ldots + \lambda v_n + \mu_1 u_1 + \ldots + \mu_m u_m$ by the
    definition of a basis. Since $v_k \in \ker T,$
    $$
    T(v) = T(\underbrace{\mu_1 u_1 + \ldots + \mu_m u_m}_{u}) = w,
    $$
    meaning there exists a $u \in U$ such that $T|_U(u) = w,$ i.e.,
    $T|_U$ is surjective.

    Additionally, assume $T|_U(u) = 0$ for some $0 \neq 0.$ Then,
    $T(\lambda_1u_1 + \ldots + \lambda_m u_m) = 0$ for some
    non-trivial $\lambda_1, \ldots, \lambda_m,$ and this vector $u
    \in \ker T.$ However, it lies outside the span of $v_1, \ldots,
    v_n,$ and thus $u \not \in \ker T.$ Thus, we derive a
    contradiction, meaning $T(u) = 0$ implies $u = 0.$ Hence, $T|_U$
    is injective.

    Since $T|_U$ is injective and surjective, $T|_U$ is an isomorphism. \qed

  \item Assume that $ST$ is invertible. Then, $ST$ is bijective. Now,
    we prove that $S, T$ are both bijective.

    Assume that $T$ is not injective, meaning there is some $v \in V
    \setminus \{0\}$ such that $Tv = 0,$ meaning $(ST)v = 0,$ meaning
    $ST$ is not injective, yielding a contradiction. By rank-nullity,
    we know $\dim V = 0 + \dim \text{range} \ T,$ and hence $T$ is
    also surjective.

    Additionally, since $ST$ is surjective, for all $w$ there exists
    some $v$ such that $STv =w ,$ meaning there exists some $u = Tv$
    such that $Su = w,$ and hence $S$ is surjective. By rank-nullity,
    $S$ is also injective.

    Therefore, both $S, T$ are bijective, meaning both $S, T$ have inverses.

    Then,m assume $S$ and $T$ are invertible. Then, $L = T^{-1}S^{-1}$ satisfies
    $$
    ST T^{-1} S^{-1} = SS^{-1} = I,
    $$
    and $$
    T^{-1} S^{-1} ST = T^{-1}T + I,
    $$
    meaning $L$ is an inverse of $ST,$ and hence $ST$ is invertible. \qed

  \item Suppose $U$ is a subspace of $V$ such that $\dim V/ U = 1.$
    Then, there is some basis of $V / U$ with a single element,
    $$\beta = \{v_1 + U\}.$$
    Hence, for all $v \in V,$ $v + U = \lambda v_1 + U,$ since
    $\beta$ is a basis. Thus, $(v - \lambda v_1) \in U,$ meaning for
    all $v \in V,$ there is some $u \in U$ such that $v = u + \lambda
    v_1$ for $u \in U.$

    First, let us assume that such $\lambda$ are unique. Then,
    construct the map $\varphi : V \to \mathbb F$
    $$
    \varphi (v) = \lambda.
    $$
    Assume $\varphi(v) = 0,$ meaning $v = u$ and hence $v \in U.$
    Thus, $\ker \varphi = U.$ Additionally, $$\varphi(v + w) =
    \varphi(u_v + \lambda_v v_1 + u_w + \lambda_w v_1) = \varphi(u_v
    + u_w + (\lambda_v + \lambda_w)v_1) = \lambda_v + \lambda_w$$
    and
    $$
    '\varphi(\mu v) = \varphi(\mu u + \mu\lambda v_1) = \mu \lambda,
    $$
    meaning $\varphi \in \mathcal L(V, \mathbb F).$

    Otherwise, there are two such $(u_1, \lambda_1), (u_2,
    \lambda_2)$ satisfying the above with $\lambda_1 \neq \lambda_2$, i.e.,
    $$
    v = u_1 + \lambda_1 v_1 = u_2 + \lambda_2 v_1.
    $$
    Then, $$v_1 = \dfrac{1}{\lambda_2 - \lambda_1}(u_1 - u_2),$$
    meaning $v_1 \in U,$ $v \in U$ and thus $V = U.$

    Construct the map $\varphi : V \to \mathbb F$ such that
    $$
    \varphi(v) = 0,
    $$
    with $\ker \varphi = V = U.$

    In either case, we can construct such a $\varphi.$ \qed

  \item First, let us prove the forward direction. Assume we have $T
    = S \circ \pi.$ Then, we know $Tv = S (v + U)$ for all $v \in V.$
    Thus, for any $u \in U,$ we have
    $$
    Tu = S(u + U) = S(U).
    $$

    Then, $T(\lambda u) = \lambda T(u) = \lambda S(U).$ But we have
    $\lambda u \in U,$a meaning $\lambda S(U) = S(U)$ for any
    $\lambda \in \mathbb F,$ forcing $S(U)= 0.$ Hence, for any $u \in
    U,$ $Tu = 0,$ meaning $u \in \ker T,$ and hence $U \subseteq \ker T.$

    For the reverse direction, assume $U \subseteq \ker T.$ Assume
    for the sake of contradiction that
    $$
    v + U = v' + u
    $$
    for some $v \neq v'.$ Then, $v - v' = U \subseteq \ker T,$
    meaning $T(v - v') = 0,$ and hence $T(v) = T(v').$

    Thus, define $S \in \mathcal L(V/U, W)$ such that
    $$
    S(v + U) = T(v),
    $$
    since it is guaranteed to be a function, yielding $T(v) = S\pi v$
    for all $v.$ Thus, if $U \subseteq \ker T,$ we have $T = S \circ
    \pi$ as required. \qed

  \item First, let us suppose that $\ker \varphi \subseteq \ker
    \beta.$ Then, for any $v \in V$ with $\varphi(v) = 0,$ we have
    $\beta(v) = 0.$ Thus, trivially, for all $v \in \ker \varphi,$
    \begin{align}
      \beta(v)  = c\varphi(v) \label{null_case}
    \end{align}
    for any constant $c \in \mathbb F.$

    % Need to prove that beta(v) varphi(w) - beta(w) varphi(v) = 0%
    Consider the term $$
    x = \varphi(w) v - \varphi(v) w \in V.
    $$
    Since $\varphi(x) = \varphi(\varphi(w)v) - \varphi(\varphi(v)w),$
    applying linearity yields
    $$
    \varphi(x) = \varphi(w) \varphi(v) - \varphi(v) \varphi(w) = 0,
    $$
    meaning $x \in \ker \varphi.$ Since $\ker \varphi \subseteq \ker
    \beta,$ $x \in \ker \beta,$ and thus
    $$
    \beta(\varphi(w)v - \varphi(v)w) = 0.
    $$
    Equivalently we have
    $$
    \varphi(w) \beta(v) = \varphi(v) \beta(w).
    $$
    When $v, w \not \in \ker \varphi,$ we know $\beta(v)/\varphi(v) =
    \beta(w)/\varphi(w)$ for all such $v, w,$ meaning
    $$
    \beta(v) = c\varphi(v)
    $$
    for some constant $c \in \mathbb F$ for all $v \not \in \ker
    \varphi.$ Combining this result with \eqref{null_case}, there
    exists some constant $c \in \mathbb F$ such that
    $$
    \beta = c \varphi
    $$
    if $\ker \varphi \subseteq \ker \beta.$

    Conversely, let us suppose that we have $\varphi, \beta \in V'$
    such that $\beta = c \varphi$ for some $c \in \mathbb F.$ Then,
    $\beta(v) = \varphi(cv)$ for all $v \in V.$ Thus, consider the
    case when $v \in \ker \varphi.$ Then, $\beta(v) = c\varphi(v) =
    c(0) = 0,$ meaning $v \in \ker \beta.$ Hence, $\ker \varphi
    \subseteq \ker \beta.$

    Therefore, $\ker \varphi \subseteq \ker \beta$ iff $\beta = c
    \varphi$ for some constant $c \in \mathbb F.$ \qed
\end{enumerate}

\subsection{Evaluating Proofs}
This proof requires $\dim V$ to be well-defined, which fails when $V$
is infinite dimensional.

\subsection{Using Computers}
\begin{center}
  \includegraphics[width=0.9\linewidth]{math3406/Screenshot
  2026-10-17 at 2.31.25 pm.png}
\end{center}

Conjecture: $\dim \mathcal E = (\dim V - \dim U) \cdot \dim W,$ since
the constraint forces $T$ to send the basis elements of $U$ to $0,$
but the remaining are free.
\end{document}
